\section{Study texts and liturgical locators}
The Scripture below is the public-domain Douay--Rheims Challoner study translation, as transmitted in the registered Gutenberg delivery. It is \textbf{not the approved English proclaimed at Mass}; the U.S. Lectionary controls that text. The spelling and wording are those of the historical witness. Verse numbers mark the biblical sources. The three readings follow today's appointed cuts; contextual Scripture is identified separately.

\subsection{Entrance and Collect}
\label{proper-entrance}
\textbf{Entrance:} \textit{Salus populi ego sum}. Roman Missal, Twenty-fifth Sunday in Ordinary Time; ICEL Antiphonary, printed p. 79. The composed antiphon is not reproduced here. Psalm 37:39--40 (Douay--Rheims 36:39--40) is its identified underlying scriptural basis, not the full antiphon's text.

\label{proper-collect}
\textbf{Collect:} \textit{Deus, qui sacrae legis}. Roman Missal, same Sunday, Collect. The complete prayer is not reproduced; it names a law of love, asks for obedience, and looks to an eternal end.

\subsection{Isaiah 55:6--9}
\label{proper-first-reading}

\textsuperscript{6} Seek ye the Lord, while he may be found: call upon him, while he is near.

\textsuperscript{7} Let the wicked forsake his way, and the unjust man his thoughts, and let him return to the Lord, and he will have mercy on him, and to our God: for he is bountiful to forgive.

\textsuperscript{8} For my thoughts are not your thoughts: nor your ways my ways, saith the Lord.

\textsuperscript{9} For as the heavens are exalted above the earth, so are my ways exalted above your ways, and my thoughts above your thoughts.

\subsection{Psalm 145:2--3, 8--9, 17--18}
\label{proper-responsorial-psalm}
Douay--Rheims Psalm 144. The appointed response is drawn from 18a; its approved liturgical wording is not reproduced. The three strophes follow.

\textsuperscript{2} Every day will I bless thee: and I will praise thy name for ever; yea, for ever and ever.

\textsuperscript{3} Great is the Lord, and greatly to be praised: and of his greatness there is no end.

\medskip

\textsuperscript{8} The Lord is gracious and merciful: patient and plenteous in mercy.

\textsuperscript{9} The Lord is sweet to all: and his tender mercies are over all his works.

\medskip

\textsuperscript{17} The Lord is just in all his ways: and holy in all his works.

\textsuperscript{18} The Lord is nigh unto all them that call upon him: to all that call upon him in truth.

\subsection{Philippians 1:20c--24, 27a}
\label{proper-second-reading}
The opening is an excerpt within verse 20; verses 25--26 and the continuation of 27 are omitted. The terminal colon in the historical wording of 27a marks the cut before its unappointed continuation.

\textsuperscript{20} [\ldots] shall Christ be magnified in my body, whether it be by life or by death.

\textsuperscript{21} For to me, to live is Christ: and to die is gain.

\textsuperscript{22} And if to live in the flesh: this is to me the fruit of labour. And what I shall choose I know not.

\textsuperscript{23} But I am straitened between two: having a desire to be dissolved and to be with Christ, a thing by far the better.

\textsuperscript{24} But to abide still in the flesh is needful for you.

\textsuperscript{27} Only let your conversation be worthy of the gospel of Christ:

\subsection{Gospel acclamation: Cf. Acts 16:14b}
\label{proper-acclamation}
The appointed acclamation is an adaptation addressed to the Lord and is not reproduced. Its underlying narrative is Acts 16:14, printed here solely as contextual Scripture:

\textsuperscript{14} And a certain woman named Lydia, a seller of purple, of the city of Thyatira, one that worshipped God, did hear: whose heart the Lord opened to attend to those things which were said by Paul.

\subsection{Matthew 20:1--16a}
\label{proper-gospel}
The appointed reading stops after the first--last reversal; the many-called, few-chosen clause that follows in the historical witness lies outside it.

\textsuperscript{1} The kingdom of heaven is like to an householder, who went out early in the morning to hire labourers into his vineyard.

\textsuperscript{2} And having agreed with the labourers for a penny a day, he sent them into his vineyard.

\textsuperscript{3} And going out about the third hour, he saw others standing in the marketplace idle.

\textsuperscript{4} And he said to them: Go you also into my vineyard, and I will give you what shall be just.

\textsuperscript{5} And they went their way. And again he went out about the sixth and the ninth hour, and did in like manner.

\textsuperscript{6} But about the eleventh hour he went out and found others standing, and he saith to them: Why stand you here all the day idle?

\textsuperscript{7} They say to him: Because no man hath hired us. He saith to them: Go ye also into my vineyard.

\textsuperscript{8} And when evening was come, the lord of the vineyard saith to his steward: Call the labourers and pay them their hire, beginning from the last even to the first.

\textsuperscript{9} When therefore they were come that came about the eleventh hour, they received every man a penny.

\textsuperscript{10} But when the first also came, they thought that they should receive more: And they also received every man a penny.

\textsuperscript{11} And receiving it they murmured against the master of the house,

\textsuperscript{12} Saying: These last have worked but one hour, and thou hast made them equal to us, that have borne the burden of the day and the heats.

\textsuperscript{13} But he answering said to one of them: friend, I do thee no wrong: didst thou not agree with me for a penny?

\textsuperscript{14} Take what is thine, and go thy way: I will also give to this last even as to thee.

\textsuperscript{15} Or, is it not lawful for me to do what I will? Is thy eye evil, because I am good?

\textsuperscript{16} So shall the last be first and the first last.

\subsection{Prayer over the Offerings}
\label{proper-offerings}
\textit{Munera, quaesumus, Domine}. Roman Missal, Twenty-fifth Sunday in Ordinary Time, Prayer over the Offerings. The complete prayer is not reproduced. It relates the people's offerings and faith to sacramental participation.

\subsection{Communion: Psalm 119:4--5, or John 10:14}
\label{proper-communion-ps119}
\textbf{First alternative: Psalm 119:4--5} (Douay--Rheims 118). Historical Scripture study text:

\textsuperscript{4} Thou hast commanded thy commandments to be kept most diligently.

\textsuperscript{5} O! that my ways may be directed to keep thy justifications.

\Needspace{4\baselineskip}
\label{proper-communion-jn10}
\textbf{Or, second alternative: John 10:14.} Historical Scripture study text:

\textsuperscript{14} I am the good shepherd: and I know mine, and mine know me.

\subsection{Prayer after Communion}
\label{proper-after-communion}
\textit{Quos tuis, Domine}. Roman Missal, Twenty-fifth Sunday in Ordinary Time, Prayer after Communion. The complete prayer is not reproduced. It asks continuing help so that redemption bears fruit in sacramental life and conduct.
